\begin{tikzpicture}[
    node distance=0.27in and 0.48in,
    function/.style={draw=black!60, fill=black!3, rounded corners=2pt,
        minimum width=2.18in, minimum height=0.28in, align=center,
        font=\ttfamily\scriptsize},
    entry/.style={function, fill=green!8, draw=green!40!black,
        line width=0.65pt},
    wall/.style={function, fill=violet!7, draw=violet!55!black},
    constraint/.style={function, fill=orange!10, draw=orange!60!black},
    force/.style={function, fill=blue!7, draw=blue!45!black},
    flow/.style={-{Stealth[length=5pt]}, draw=black!65, line width=0.45pt},
    call/.style={-{Stealth[length=4pt]}, draw=blue!55!black, dashed,
        line width=0.4pt}
]

\newcommand{\ForceFlowNode}[2]{\hyperref[force:#1]{\strut #2}}

\node[entry] (main) {per-source compute invocation};

\node[wall, below left=0.46in and 1.42in of main] (piston)
    {\ForceFlowNode{ProcessPistonCollision}{ProcessPistonCollision}};
\node[function, below=0.46in of main] (physical)
    {\ForceFlowNode{GetPhysicalParticleContact}{GetPhysicalParticleContact}};
\node[wall, below right=0.46in and 1.42in of main] (marker)
    {\ForceFlowNode{ParametricBoundaryMarkerApplies}{ParametricBoundaryMarkerApplies}};

\node[wall, below=of piston] (piston-eval)
    {\ForceFlowNode{EvaluatePistonWall}{EvaluatePistonWall}};
\node[function, below=of physical] (particle)
    {\ForceFlowNode{ProcessParticleCollision}{ProcessParticleCollision}};
\node[wall, below=of marker] (wall-contacts)
    {\ForceFlowNode{EvaluateParametricWallContacts}{EvaluateParametricWallContacts}};

\node[function, below=of particle] (particle-state)
    {\ForceFlowNode{InitializeContactState}{InitializeContactState}};
\node[wall, below=of wall-contacts] (wall-eval)
    {\ForceFlowNode{EvaluateParametricWallSegment}{EvaluateParametricWallSegment}};
\node[wall, below=of wall-eval] (wall-process)
    {\ForceFlowNode{ProcessParametricWallCollision}{ProcessParametricWallCollision}};
\node[wall, below=of wall-process] (wall-state)
    {\ForceFlowNode{InitializeParametricWallContactState}{InitializeParametricWallContactState}};

\node[force, below=2.00in of particle-state] (force)
    {\ForceFlowNode{AccumulateContactForce}{AccumulateContactForce}};
\node[entry, below=of force] (velocity)
    {\ForceFlowNode{CalcVelocity}{CalcVelocity}};
\node[constraint, below=of velocity] (maximum)
    {\ForceFlowNode{ApplySourceMaximumDepth}{ApplySourceMaximumDepth}};
\node[constraint, below=of maximum] (project)
    {\ForceFlowNode{ProjectSourceVelocityToContactSet}{ProjectSourceVelocityToContactSet}};
\node[entry, below=of project] (position)
    {\ForceFlowNode{CalcPosition}{CalcPosition}};

\draw[flow] (main) -- (piston);
\draw[flow] (main) -- (physical);
\draw[flow] (main) -- (marker);

\draw[flow] (piston) -- (piston-eval);
\draw[flow] (physical) -- (particle);
\draw[flow] (particle) -- (particle-state);
\draw[flow] (marker) -- (wall-contacts);
\draw[flow] (wall-contacts) -- (wall-eval);
\draw[flow] (wall-eval) -- (wall-process);
\draw[flow] (wall-process) -- (wall-state);

\draw[call] (piston-eval.south) |- (force.west);
\draw[call] (particle-state.south) -- (force.north);
\draw[call] (wall-state.south) |- (force.east);
\draw[flow] (force) -- (velocity);
\draw[flow] (velocity) -- (maximum);
\draw[flow] (maximum) -- (project);
\draw[flow] (project) -- (position);

\end{tikzpicture}
